\documentclass[10pt,a4paper]{amsart}
\usepackage[margin=19mm]{geometry}
\usepackage{amsmath,amssymb,mathtools}
\usepackage[scr=rsfs,cal=euler]{mathalfa}
\usepackage{xfrac}
\usepackage[unicode,hidelinks,bookmarks=false]{hyperref}
\newcommand{\ie}{i.e.\ }
\newcommand{\cf}{cf.\ }
\newcommand{\resp}{respectively\ }
\newcommand{\Subsection}{\subsection}
\providecommand{\nobreakdash}{\nobreak\hbox{-}}
\setlength{\emergencystretch}{2em}
\multlinegap0pt
\usepackage{bm}
\usepackage{amssymb}
\usepackage{algorithm}\usepackage{algpseudocode}
\usepackage{xcolor}
\usepackage[shortlabels]{enumitem}
\newtheorem{theorem}{Theorem}
\newtheorem{corollary}{Corollary}
\title{Spectral Approximation and Ergodic-Capacity Convergence of HMIMO Channels under Spatial–Wavenumber Domain Mismatch}
\author{Hangsong Yan, \textit{Member}, \textit{IEEE}, Hong Yang, \textit{Senior Member}, \textit{IEEE}, Shu Sun, \textit{Senior Member}, \textit{IEEE}}
\date{Source: arXiv:2608.26802v1 (CC BY 4.0)}
\begin{document}
\maketitle

\noindent\emph{Source and licence.} Spectral Approximation and Ergodic-Capacity Convergence of HMIMO Channels under Spatial–Wavenumber Domain Mismatch. Version arXiv:2608.26802v1 (CC BY 4.0), distributed by arXiv under the Creative Commons Attribution 4.0 International licence, http://creativecommons.org/licenses/by/4.0/, as recorded in the arXiv licence metadata for this exact version and shown on the abstract page https://arxiv.org/abs/2608.26802v1. Full licence notice: \texttt{licenses/arxiv-2608.26802-CC-BY-4.0.txt}.\par\smallskip

\noindent\textit{Editorial scope.} Counted unit: the article's corollary corollary:eigenspectrum\_decay (source0.tex lines 479--503). The article's complete original proof of it is delivered with it (source0.tex lines 504--519, one \texttt{proof} environment of 16 lines). No mathematics is written, repaired or re-derived: the statement and the proof are the article's own lines, in the article's own notation, and no retained line is substituted or re-flowed.  Retained uncounted as the source-local context the proof needs: lines 411--454.  Nothing was omitted from any retained span, so the package carries no omission record.  No retained line contains a citation, so the wrapper carries no bibliography and no bibliography entry.  No retained line contains a figure environment, an image-inclusion command or drawing code, so no image is delivered and the package declares no asset dependency.  Editorial formatting only, in addition: an AMS \texttt{amsart} preamble that restates the article's own declarations and macros used by the retained lines, namely the environments \texttt{theorem}, \texttt{corollary} on separate counters, together with the standard packages those lines need.  The retained environments print in this document as Theorem 1, Corollary 1 in order of appearance, whereas the article prints the same items with the numbering of its own full document.  The complete native source of the article as distributed in the arXiv e-print for this exact version is bundled unchanged as \texttt{sources/208661/source0.tex}.
\begin{theorem}
	\label{theorem:approx_error}
	Let $\mu_n(\mathcal{P})$ be the $n$-th eigenvalue of the
	non-separable continuous spatial-wavenumber integral operator
	$\mathcal{P}$, arranged in non-increasing order, and let
	$\mu_n(\mathbf{M}_N)$ be the $n$-th eigenvalue of the
	finite-dimensional truncated matrix $\mathbf{M}_N$, also arranged
	in non-increasing order, with the symmetric truncation
	$N=N_{1D}^2$. If
	$$
		N_{1D}>\frac{ec}{4},
	$$
	then the total absolute spectral approximation error
	\begin{align}
		&\sum_{n=0}^{N-1}
		\left|
		\mu_n(\mathcal{P})-\mu_n(\mathbf{M}_N)
		\right|
		+
		\sum_{n=N}^{\infty}\mu_n(\mathcal{P})
		\nonumber\\
		= & 
		\operatorname{Tr}(\mathcal{P})
		-
		\operatorname{Tr}(\mathbf{M}_N) <
		C_1(N_{1D},c)
		e^{-2N_{1D}
			\ln\left(\frac{4N_{1D}}{ce}\right)}, \nonumber
		\label{eq:error_truncation}
	\end{align}
	where
	\begin{equation}
		C_1(N_{1D},c)
		=
		\frac{c^2}{\pi N_{1D}}
		\left[
		1-
		\left(
		\frac{ec}{4(N_{1D}+1)}
		\right)^2
		\right]^{-1}.
		\label{eq:C1_definition}
	\end{equation}
\end{theorem}

\begin{corollary}
	\label{corollary:eigenspectrum_decay}
	Define $m_n \triangleq \left\lfloor\sqrt{n}\right\rfloor.$
	If $m_n>ec/4$, then the $n$-th eigenvalue of the continuous
	operator $\mathcal{P}$ satisfies
	\begin{align}
		\mu_n(\mathcal{P})
		<
		\frac{C_1(m_n,c)}
		{n-m_n^2+1}
		e^{-2m_n
			\ln\left(\frac{4m_n}{ce}\right)}.
		\label{eq:UB_eigenvalue}
	\end{align}
	Consequently, the flattened two-dimensional eigenspectrum
	admits the asymptotic upper envelope
	\begin{equation}
		\ln\mu_n(\mathcal{P})
		\leq
		-2\sqrt{n}
		\ln\left(\frac{4\sqrt{n}}{ce}\right)
		+\mathcal{O}(\ln n).
		\label{eq:flattened_asymptotic}
	\end{equation}
\end{corollary}

\begin{proof}
	Since $m_n^2\leq n<(m_n+1)^2$ and the eigenvalues of
	$\mathcal{P}$ are arranged in non-increasing order,
	$$
		(n-m_n^2+1)\mu_n(\mathcal{P})
		\leq
		\sum_{k=m_n^2}^{n}\mu_k(\mathcal{P})\leq
		\sum_{k=m_n^2}^{\infty}\mu_k(\mathcal{P}).
	$$
	Applying Theorem~\ref{theorem:approx_error} with the
	one-dimensional truncation dimension set to $m_n$ yields
	\eqref{eq:UB_eigenvalue}. Since
	$m_n=\sqrt{n}+\mathcal{O}(1)$, taking the logarithm of
	\eqref{eq:UB_eigenvalue} further gives
	\eqref{eq:flattened_asymptotic}.
\end{proof}

\end{document}
